%! Unit 1 cover — page 1. THE SHEET ITSELF LIVES HERE.
%! Both unit_cover/main.tex and unit_cover_key/main.tex \input this file, so
%! page 1 cannot drift between the student packet and the key packet. Edit the
%! cover here, never in a wrapper.

% Full-bleed cerulean banner
\begin{tikzpicture}[remember picture, overlay]
  \fill[cerulean] (current page.north west)
    rectangle ([yshift=-1.25in]current page.north east);
\end{tikzpicture}
\vspace{-0.72in}

\begin{center}
  {\color{white}\textbf{\LARGE Statistical Analysis \& Algebraic Reasoning}}\\[4pt]
  {\color{frostmid}\large\textbf{Unit 1}}\\[6pt]
  {\color{white}\textbf{\Large Asking Questions with Data: Study Design}}\\[5pt]
  {\color{frostmid}\small\textbf{Standards: PS.DC.1 \quad PS.DC.2 \quad PS.DC.3 \quad AFDA.DA.2}}
\end{center}

\vspace{0.20in}
\namedateperiod
\vspace{0.14in}

\begin{objectivebox}
\small
Every number you will ever read about people came from a study somebody
designed. This unit is about \textbf{where data comes from} --- how a question
gets asked, who gets asked, and what the answer is then allowed to claim. By
the end you can plan a study, name what is wrong with a bad one, and tell the
difference between two sentences that look identical: \emph{is associated with}
and \emph{caused}.
\end{objectivebox}

\vspace{0.12in}

\begin{tocbox}
\small
\renewcommand{\arraystretch}{1.22}
\begin{tabularx}{\linewidth}{@{} c X l @{}}
\rowcolor{frost}
\textbf{\#} & \textbf{Lesson} & \textbf{Standards} \\
\midrule
1.0 & Unit Launch: Study Design                        & PS.DC.1a--c; AFDA.DA.2a \\
1.1 & The Statistical Cycle and Types of Data          & PS.DC.1a--c; AFDA.DA.2a \\
1.2 & Populations, Samples, Parameters, and Statistics & PS.DC.1d--e \\
1.3 & Choosing a Sample --- Four Sampling Techniques   & PS.DC.2a--b; AFDA.DA.2b \\
1.4 & Bias in Samples and Surveys                      & PS.DC.2c; AFDA.DA.2e--g \\
1.5 & Observational Studies                            & PS.DC.2d \\
1.6 & Principles of Experimental Design                & PS.DC.3a--b; AFDA.DA.2c \\
1.7 & Comparing Studies and Choosing a Method          & PS.DC.3c--e \\
1.8 & Project: Design and Conduct a Survey             & AFDA.DA.2d, h--j; PS.DC.2d \\
\midrule
    & \multicolumn{1}{r}{\textbf{Practice Test}} & \textbf{all of Unit 1} \\
\end{tabularx}
\end{tocbox}

\vspace{0.12in}

\begin{spiralbox}
\small
\begin{enumerate}[leftmargin=1.6em, itemsep=3pt, topsep=2pt]
  \item \textbf{The question decides the data.} What you ask determines whether
        the answers are categories or numbers --- before a single value exists.
  \item \textbf{A sample speaks only for the group chance could have reached.}
        A statistic describes the sample; a parameter describes the population.
  \item \textbf{Bias is a direction, not bad luck.} Honest samples scatter
        around the truth; a biased method misses the same way every time --- and
        a bigger sample repeats the error more precisely.
  \item \textbf{Who decides the groups decides the verb.} If people sorted
        themselves, you may write \emph{is associated with}. Only random
        assignment earns \emph{caused}.
\end{enumerate}
\end{spiralbox}

\vspace{0.12in}

\begin{remindbox}
\small
\textbf{The practice test at the back of this packet is yours to keep.} It
matches the real Unit 1 test in length, format, and ideas --- same four parts,
same $100$ points --- with different numbers and contexts. Work it with no
notes before you check the key, and pay attention to the questions that ask for
a \emph{direction} or a \emph{justification}: on this unit's test, naming the
right term without saying which way the error runs, or why the design earns the
verb, is worth only half the points.
\end{remindbox}
